% Cote 119: offprints, reports, courses, CV and letters (typed and printed part): archive cote 119, site page 138 (PDF page 139).
% typescript. Draft compared: none.
% AI-transcribed by Claude Opus 5.5 (claude-opus-5-5, adaptive thinking, effort high), Message Batches API, batch r119_full (msgbatch_01DwZNm5Y7YsKBWyAN4sFs8w), 2026-11-13; not yet checked by eye.
\scanpage{119/138}
notions et idées autour de la notion de \uline{motif}\hand{,} que j'ai introduites\inserted[typed]{et développées} tout
au long des années soixante (tombées dans un oubli soudain dès mon \xout{départ}
"départ" en 1970 et -- à une exhumation p\unsure{\hand{a}}rtielle près en 1982 -- jus-%
qu'à aujourd'hui même ...), ainsi qu'une maîtrise des principales \xout{idées}
notions et idées de \uline{géométrie algébrique anabélienne} que j'ai dégagées
depuis 1978. En termes imagés, on peut dire que ces deux courants d'
idées, le courant "abélien" incarné par la notion de motif, et le
courant "anabélien" exemplifié par la structure géométrico-arithmétique
de la "tour de Teichmüller", sont à la "géométrie arithmétique" dans son
enfance, ce que les \xout{\unsure{notions}} courants "complexes de cochaines -- catégories
dérivées commutatives" et "champs en tous genres -- catégories homotopi-%
ques non commutatives" sont à l'algèbre topologique\unsure{\hand{,}} (encore in utero).

\margin{\unsure{$\lfloor$}}La différence essentielle entre ces deux disciplines ne se situe cepen-%
dant nullement dans leur état d'avancement relatif, circonstance des
plus contingentes\hand{,} apte à changer radicalement en l'espace de quelques
années, ne serait-ce que par la vertu de l'écriture de "Pursuing Stacks"\hand{.}
Elle se situe plutôt dans une différence de profondeur. Comme c'était
le cas naguère pour le développement d'une"topologie générale" (faite
sur mesure pour l'analyste, bien plus que pour le géomètre), le travail
essentiel à accomplir pour la mise sur pied de l'"algèbre topologique"
est\hand{,} avant toute autre chose, le \uline{développement d'un langage} , dont le
manque se fait sentir (à moi tout au moins) à tous les pas. La dimension
de cette tâche me semble être du même ordre que celle à laquelle se sont
vu\hand{s} confronté\hand{s} Hausdorf et d'autres, dans l'entre-deux-guerres. Elle n'a
pas de commune mesure avec la tâche posée par le développement de la
géométrie arithmétique, jusqu'au point d'un sentiment de "maîtrise" des
principaux courants d'idées qui y confluent. On mesurera cette différen-%
ce de "dimensions", en se rappelant que la "maîtrise" du\hand{"}courant
\deleted{"}anabélien" impliquerait, nota\inserted{m}ment, une description "purement algébri-%
que" du groupe de Galois de $\overline{\underline{Q}}$ sur $\underline{Q}$ , et de la famille de ses sous-%
groupes de décomposition et d'inertie associés aux différents nombres pre-%
miers. La structure de ces derniers \xout{\unsure{a été}} (correspondants aux cas
"locaux" des corps p-adiques) ont été déterminés récemment par \xout{\unsure{Uwe Je}}
Uwe Jennsen, Kay Wingberg et (dans le cas p=2) Volker Diekert, avec une
relation principale rappelant étrangement celle qui décrit le groupe fon-%
damental d'une courbe algébrique sur le corps des complexes. Mais on est
loin sans doute d'une description analogue dans le cas "global".

\mspage{138}
